\documentclass[a4,11pt]{aleph-notas}

% -- Paquetes adicionales
\usepackage{enumitem}
\usepackage{aleph-comandos}
\usepackage{systeme}

% -- Datos
\institucion{Facultad de Ciencias Exactas, Naturales y Ambientales}
\carrera{Catálogo STEM}
\asignatura{Álgebra lineal}
\tema{Ejercicios 04: Determinantes}
\autor[A. Merino]{Andrés Merino}
\fecha{Periodo 2026-1} 

\logouno[0.16\textwidth]{Logos/logoPUCE_04_ac}
\definecolor{colortext}{HTML}{1957d1}
\definecolor{colordef}{HTML}{1957d1}
\fuente{montserrat}

\begin{document}

\encabezado

\section{Determinantes}

%%%%%%%%%%%%%%%%%%%%%%%%%%%%%%%%%%%%%%%%
\begin{ejer}
    Sea la matriz:
    \[
        A= \begin{pmatrix}
            \alpha & -1 & 0 \\
            -2 & \alpha & -2 \\
            0 & -1 & \alpha
        \end{pmatrix},
    \]
    donde $\alpha \in \mathbb{R}$.
    \begin{enumerate}
    \item 
        ¿Para qu\'{e} valores de $\alpha$ la matriz $A$ es invertible? 
    \item 
        Calcule la inversa de $A$ cuando sea posible.
    \end{enumerate}
\end{ejer}

\begin{proof}[Solución]\hspace{0pt}
    \begin{enumerate}
    \item 
    Para saber si A es invertible, se debe calcular el determinante: 
    \[
    det(A)=\begin{vmatrix}
                \alpha & -1 & 0 \\
                -2 & \alpha & -2 \\
                0 & -1 & \alpha
                \end{vmatrix}=\alpha \begin{vmatrix}
                \alpha & -2 \\
                -1 & \alpha
                \end{vmatrix}-\begin{vmatrix}
                -2 & -2 \\
                0 & \alpha
                \end{vmatrix}=\alpha^3-4\alpha.\\
    \]
    Por lo tanto, $A$ es invertible si y solo si $\alpha \neq-2$, $\alpha\neq0$ y $\alpha\neq2$.
    
    \item Supongamos que $\alpha \in \mathbb{R}\setminus\{-2, 0, 2\}$. Entonces, $A^{-1}$ existe y tenemos:
    \[
    A^{-1}=\frac{1}{det(A)}Adj(A)=\frac{1}{\alpha^3-4\alpha}   \begin{pmatrix}
                \alpha^2-2 & \alpha & 2 \\
                2\alpha & \alpha^2 & 2\alpha \\
                2 & \alpha & \alpha^2-2
            \end{pmatrix}.\qedhere
    \]
    \end{enumerate}
\end{proof}

%%%%%%%%%%%%%%%%%%%%%%%%%%%%%%%%%%%%%%%%
\begin{ejer}
    Sea
    \[
        B=
        \begin{pmatrix}
            1 & 2 & -1 \\
            2 & 4 & -2 \\
            0 & 3 & 5
        \end{pmatrix}.
    \]
    
    \begin{enumerate}
    \item 
        Calcule $\det(B)$.
    \item
        Determine si $B$ es invertible.
    \end{enumerate}
\end{ejer}

\begin{proof}[Solución]\hspace{0pt}
    \begin{enumerate}
    \item
        Observamos que la segunda fila es el doble de la primera:
        \[
            F_2 = 2F_1.
        \]
        Por la propiedad de determinantes, si dos filas son proporcionales, el determinante es cero. Por tanto,
        \[
            \det(B)=0.
        \]
    \item
        Una matriz cuadrada es invertible si y solo si su determinante es distinto de cero. Como
        \[
            \det(B)=0,
        \]
        concluimos que $B$ no es invertible.\qedhere
    \end{enumerate}
\end{proof}

%%%%%%%%%%%%%%%%%%%%%%%%%%%%%%%%%%%%%%%%
\begin{ejer}
    Sea
    \[
        C=
        \begin{pmatrix}
            2 & 1 & 0 \\
            -1 & 3 & 2 \\
            4 & 0 & 1
        \end{pmatrix}.
    \]
    
    \begin{enumerate}
    \item
        Calcule $\det(C)$ usando expansión por cofactores en la primera fila.
    \item
        Halle el cofactor $C_{23}$.
    \end{enumerate}
\end{ejer}

\begin{proof}[Solución]\hspace{0pt}
    \begin{enumerate}
    \item
        Expansión por la primera fila:
        \[
            \det(C)=2\begin{vmatrix}3 & 2 \\ 0 & 1\end{vmatrix}-1\begin{vmatrix}-1 & 2 \\ 4 & 1\end{vmatrix}+0\begin{vmatrix}-1 & 3 \\ 4 & 0\end{vmatrix}.
        \]
        Entonces,
        \[
            \det(C)=2(3)-1((-1)(1)-2(4))=6-(-9)=15.
        \]
    \item
        El menor $23$ se obtiene eliminando la fila 2 y la columna 3:
        \[
            M_{23}=\begin{vmatrix}2 & 1 \\ 4 & 0\end{vmatrix}=2(0)-1(4)=-4.
        \]
        Por definición de cofactor,
        \[
            C_{23}=(-1)^{2+3}M_{23}=(-1)^5(-4)=4.\qedhere
        \]
    \end{enumerate}
\end{proof}

%%%%%%%%%%%%%%%%%%%%%%%%%%%%%%%%%%%%%%%%
\begin{ejer}
    Sean
    \[
        A=
        \begin{pmatrix}
            1 & 2 \\
            3 & 1
        \end{pmatrix},
        \quad
        D=
        \begin{pmatrix}
            2 & -1 \\
            0 & 4
        \end{pmatrix}.
    \]
    
    \begin{enumerate}
    \item
        Calcule $\det(A)$ y $\det(D)$.
    \item
        Verifique que $\det(AD)=\det(A)\det(D)$.
    \item
        Calcule $\det(3A)$ y compare con $3^2\det(A)$.
    \end{enumerate}
\end{ejer}

\begin{proof}[Solución]\hspace{0pt}
    \begin{enumerate}
    \item
        \[
            \det(A)=1\cdot1-2\cdot3=-5,
            \qquad
            \det(D)=2\cdot4-(-1)\cdot0=8.
        \]
    \item
        Primero,
        \[
            AD=
            \begin{pmatrix}
                1 & 2 \\
                3 & 1
            \end{pmatrix}
            \begin{pmatrix}
                2 & -1 \\
                0 & 4
            \end{pmatrix}
            =
            \begin{pmatrix}
                2 & 7 \\
                6 & 1
            \end{pmatrix}.
        \]
        Luego,
        \[
            \det(AD)=2\cdot1-7\cdot6=-40.
        \]
        Además,
        \[
            \det(A)\det(D)=(-5)(8)=-40.
        \]
        Por tanto, se verifica que
        \[
            \det(AD)=\det(A)\det(D).
        \]
    \item Calculemos:
        \[
            3A=
            \begin{pmatrix}
                3 & 6 \\
                9 & 3
            \end{pmatrix},
            \qquad
            \det(3A)=3\cdot3-6\cdot9=-45.
        \]
        Por otro lado,
        \[
            3^2\det(A)=9(-5)=-45.
        \]
        Se cumple la propiedad $\det(\alpha A)=\alpha^n\det(A)$ con $n=2$.\qedhere
    \end{enumerate}
\end{proof}

%%%%%%%%%%%%%%%%%%%%%%%%%%%%%%%%%%%%%%%%
\begin{ejer}
    Sea
    \[
        E=
        \begin{pmatrix}
            2 & 1 & 1 \\
            1 & 1 & 0 \\
            1 & 0 & 2
        \end{pmatrix}.
    \]
    
    \begin{enumerate}
    \item
        Calcule $\det(E)$.
    \item
        Obtenga $E^{-1}$ usando la matriz adjunta.
    \item
        Resuelva el sistema $Ex=b$, con $b=(1,2,3)^\intercal$.
    \end{enumerate}
\end{ejer}

\begin{proof}[Solución]\hspace{0pt}
    \begin{enumerate}
    \item
        Expansión por la primera fila:
        \[
            \det(E)=2\begin{vmatrix}1 & 0 \\ 0 & 2\end{vmatrix}-1\begin{vmatrix}1 & 0 \\ 1 & 2\end{vmatrix}+1\begin{vmatrix}1 & 1 \\ 1 & 0\end{vmatrix}=2(2)-1(2)+1(-1)=1.
        \]
        Así, $\det(E)=1$.
    \item
        Calculamos cofactores:
        \[
            \cof(E)=
            \begin{pmatrix}
                2 & -2 & -1 \\
                -2 & 3 & 1 \\
                -1 & 1 & 1
            \end{pmatrix}.
        \]
        Luego,
        \[
            \adj(E)=\cof(E)^\intercal=
            \begin{pmatrix}
                2 & -2 & -1 \\
                -2 & 3 & 1 \\
                -1 & 1 & 1
            \end{pmatrix}.
        \]
        Como $\det(E)=1$,
        \[
            E^{-1}=\frac{1}{\det(E)}\adj(E)=
            \begin{pmatrix}
                2 & -2 & -1 \\
                -2 & 3 & 1 \\
                -1 & 1 & 1
            \end{pmatrix}.
        \]
    \item
        La solución es $x=E^{-1}b$:
        \[
            x=
            \begin{pmatrix}
                2 & -2 & -1 \\
                -2 & 3 & 1 \\
                -1 & 1 & 1
            \end{pmatrix}
            \begin{pmatrix}
                1 \\
                2 \\
                3
            \end{pmatrix}
            =
            \begin{pmatrix}
                -5 \\
                7 \\
                4
            \end{pmatrix}.
        \]
        Por tanto, la solución del sistema es
        \[
            x_1=-5,\quad x_2=7,\quad x_3=4.\qedhere
        \]
    \end{enumerate}
\end{proof}



\end{document}